\section{Conception}
\label{sec:conception}

\subsection{Overview}
\label{sec:concept-overview}

The concept has two connected parts. The first part generates a scene graph
of the pick-and-place work cell from the robot's cameras and the robot state.
The second part is an agentic framework that inspects this graph and the
camera images with tools and annotates the graph with risk-relevant
information: object states such as a missing cork, label content such as a
corrosive pictogram, object parts such as the blade of a knife, and the
relations that turn these details into hazards. The annotated graph is
produced while the task runs, but it does not control the robot.

The pipeline consists of five modules connected through ROS, as shown in
Figure~\ref{fig:architecture}. A sensor layer provides time-aligned camera
images, depth data and robot joint states, together with the camera poses
from calibration. A perception module detects objects and object parts and
estimates their 3D poses. A scene graph generator builds and updates the
graph from these results. An annotation agent inspects the graph and the
images and writes risk annotations. Finally, a validation and export module
checks every annotation against the schema and publishes the annotated
graph. Because all modules read their input from ROS topics, the pipeline
runs the same way on the live system and on recorded ROS bags. Code,
schemas, prompts and model settings are kept under version control, so
every result can be reproduced.

\begin{figure}[htbp]
  \centering
  \resizebox{\textwidth}{!}{%
  \begin{tikzpicture}[
      node distance = 6mm and 7mm,
      module/.style = {draw, rounded corners=2pt, thick, align=center,
                       minimum height=16mm, text width=24mm,
                       font=\small},
      io/.style     = {draw, dashed, align=center, minimum height=10mm,
                       text width=36mm, font=\small},
      flow/.style   = {-{Stealth[length=2.5mm]}, thick},
      aux/.style    = {-{Stealth[length=2.5mm]}, thick, dashed}
    ]
    \node[module] (sensor) {Sensor\\layer};
    \node[module, right=of sensor] (percep) {Perception};
    \node[module, right=of percep] (graph)  {Scene graph\\generator};
    \node[module, right=of graph]  (agent)  {Annotation\\agent};
    \node[module, right=of agent]  (export) {Validation\\and export};

    \node[io, above=of sensor] (input)
      {Two RGB-D cameras,\\joint states\\(live or ROS bag)};
    \node[io, below=of export] (output)
      {Annotated scene graph\\(ROS, JSON, Markdown)};

    \draw[flow] (input)  -- (sensor);
    \draw[flow] (sensor) -- (percep);
    \draw[flow] (percep) -- (graph);
    \draw[flow] (graph)  -- (agent);
    \draw[flow] (agent)  -- (export);
    \draw[flow] (export) -- (output);

    \draw[aux] (sensor.south) -- ++(0,-8mm) -| (agent.south)
      node[pos=0.25, below, font=\footnotesize] {camera images};
  \end{tikzpicture}%
  }
  \caption{Architecture of the pipeline. Solid arrows show the main data
    flow; the dashed arrow shows that the annotation agent also inspects
    the camera images.}
  \label{fig:architecture}
\end{figure}

\subsection{System Setup}
\label{sec:system-setup}

The work cell is a tabletop with a robot arm (a Franka Emika arm or another
available industrial arm), its gripper and RGB-D cameras, integrated through
ROS. Two Intel RealSense D455 cameras are used: one on the wrist for
close-up views and one fixed in front of the workspace for a third-person
view. Because the base of the arm is fixed, the camera poses follow from
the robot's kinematics and calibration, so no SLAM or visual odometry is
needed.

The concept uses existing perception models and, where suitable, the
institute's existing risk-annotated scene graph framework. The agentic
framework is built on existing LLMs and VLMs. Training models from scratch,
new motion planning or control methods, and real hazardous scenarios are
out of scope.

\subsection{Scene Graph Generation}
\label{sec:scene-graph-generation}

The scene graph describes the work cell as objects with attributes (nodes)
and relations between them (edges), for example \emph{vial in rack} or
\emph{knife held by gripper}. A versioned schema defines the node and edge
types, spatial and contact relations, persistent IDs, 3D poses, timestamps,
sources and confidence values. Nodes cover the robot, objects, object
parts, object states and labels.

Existing models detect and segment the objects and relevant object parts,
and their 3D poses are estimated from the depth data and camera
calibration. Spatial relations are computed from the 3D poses with defined
rules and thresholds. The graph is updated at a fixed rate, object
identities are kept across frames and occlusions, and each graph version is
linked to the current phase of the pick-and-place task. The graph is
published through ROS and stored as JSON, and it can be regenerated from
recorded ROS bags.

\subsection{Risk Annotation}
\label{sec:risk-annotation}

Risk information is kept in a separate annotation schema that refers to the
scene graph nodes. It supports unknown and novel values. A hazard is
expressed as a triple, for example
\begin{center}
  \texttt{(knife\_03.blade, pointing\_toward, human\_hand\_zone)}
\end{center}
with hazard type, severity, evidence and confidence attached, and the
triples are exported as BRAAIS risk triples.

The agentic framework first identifies the nodes and relations that may be
relevant for risk. It then uses vision-language models and supporting tools
to inspect objects, analyse spatial relationships and retrieve relevant
knowledge, and it creates or updates the annotations on the graph. All tool
calls and intermediate results are logged.

For each annotation, the framework records the supporting evidence,
timestamps, model details, prompts and confidence. Annotations without
valid evidence are rejected by schema validation, and uncertain or
contradictory observations are marked as unknown. The annotated graph is
published through ROS and exported as JSON or Markdown.

\subsection{Evaluation Concept}
\label{sec:evaluation-concept}

The framework is evaluated on three controlled pick-and-place scenarios,
with the chemical vial and cutlery scenarios reused from the BRAAIS
specification. It is compared against three baselines: an unannotated
scene graph, a single VLM call without tools, and a video-based VLM
approach without a scene graph. The evaluation checks whether the scene
graph is correct, whether the risk-relevant facts are correct, whether the
annotations are supported by evidence, and whether the system is fast and
reliable enough to run during the task.

For each recorded scene, reference graphs and reference annotations are
created following an annotation manual that defines how objects, parts,
relations and risk facts are labelled. The references are reviewed and
validated against ground truth, and any discrepancies are documented. At
least two LLM or VLM back-ends are tested for the agent, at least one of
them hosted locally, with repeated runs, and the effect of the main design
choices, such as the use of tools, abstention and the number of camera
views, is analysed.
